% STATUS NOTICE (2026-07-31): CANDIDATE COMPLETION, NOT YET CANONICAL.
% This note reopens the compact-psi profile route after the torsion-free
% single-section no-go.  Promoting it requires an author-approved amendment of
% AXIOM C, because the present canonical lock forbids compact-fiber averaging.

\documentclass[11pt,a4paper]{article}
\usepackage[margin=2.35cm]{geometry}
\usepackage{amsmath,amssymb,amsthm,mathtools,mathrsfs,array,tabularx}
\usepackage[most]{tcolorbox}
\usepackage{hyperref}

\newtheorem{definition}{Definition}[section]
\newtheorem{theorem}[definition]{Theorem}
\newtheorem{proposition}[definition]{Proposition}
\newtheorem{corollary}[definition]{Corollary}
\theoremstyle{remark}
\newtheorem{remark}[definition]{Remark}

\newcommand{\B}{\mathbb B}
\newcommand{\Nzero}{\mathcal N_0}
\newcommand{\avg}[1]{\left\langle #1\right\rangle_\psi}
\newcommand{\dd}{\mathrm d}
\newcommand{\one}{\mathbf 1}

\title{Candidate Local GR Completion from the Existing Complex-Time Profile\\
\large Free-embedding lift, pure-$\Theta$ Gauss action, and exact vacuum closure}
\author{Ing. David Jaro\v{s}}
\date{31 July 2026}

\begin{document}
\maketitle

\begin{abstract}
The torsion-free single-section construction
$E_\mu=\Nzero^{-1/2}D_\mu\Theta\in W_L$ obeys a concurrent-vector no-go and
cannot generate a generic non-flat vacuum geometry.  This note proves a
precise alternative using no new fundamental field: retain the already
postulated periodic complex-time profile $\psi\mapsto\Theta(x,\psi)$ and read
the four-dimensional metric from the normalized Haar average of the central
Lorentz product.  A local free isometric embedding of any Lorentzian
four-metric into $\mathbb E^{13,1}$ lifts isometrically to fourteen orthonormal
$W_L$-valued profile modes of the same $\Theta$.  The ten normal
second-derivative profiles are then independent, so variation of the
pure-$\Theta$ Gauss action is exactly equivalent to the vacuum
Einstein--$\Lambda$ equations on that branch.  This closes local
representability and local vacuum closure for the profile metric.  It does not
derive $\kappa$, $\Lambda$, profile selection, or the profile readout from the
previous canonical master action; promotion therefore requires a deliberate
canonical decision rather than a silent status change.
\end{abstract}

\begin{tcolorbox}[colback=yellow!5,colframe=yellow!55!black,title=Canonical impact]
The construction below is mathematically compatible with one fundamental
field $\Theta(q,\tau)$ and adds no independent connection or tetrad.  It is
nevertheless \emph{not compatible with the current wording of AXIOM C}, which
forbids compact-$\psi$ averaging as the canonical metric readout.  The result
is therefore a candidate completion until the repository maintainer either
promotes it by a major axiom revision or rejects the profile readout.
\end{tcolorbox}

\section{Central profile pairing}

Let
\[
 W_L=\{ix^0\one+x^k\mathbf e_k\mid x^a\in\mathbb R\}\subset\B
\]
be the canonical Lorentz slice.  For $U,V:S^1_{R_\psi}\to W_L$ define the
real symmetric profile pairing by the central identity
\begin{equation}
 \avg{U,V}\,\one
 :=\frac1{L_\psi}\int_0^{L_\psi}
 \frac12\bigl(U^\sharp V+V^\sharp U\bigr)\,\dd\psi,
 \qquad L_\psi=2\pi R_\psi.
 \label{eq:profilepair}
\end{equation}
It is the orthogonal direct sum of countably many copies of the Lorentz form
on $W_L$ and therefore has infinitely many positive and negative directions.

\begin{definition}[Candidate profile metric]
For a smooth periodic field $\Theta:U\times S^1_{R_\psi}\to W_L$ set
\begin{equation}
 E_\mu:=\partial_\mu\Theta,
 \qquad
 \boxed{\quad g_{\mu\nu}[\Theta]
 =\frac1{\Nzero}\avg{E_\mu,E_\nu}.\quad}
 \label{eq:profilemetric}
\end{equation}
Equivalently, $g_{\mu\nu}\one$ is the normalized Haar average of the
pointwise central anticommutator.  No trace or real-part projection is used.
\end{definition}

\begin{proposition}[Finite Lorentz ambient block inside one profile]
There exist fourteen smooth periodic profiles $F_A:S^1\to W_L$,
$A=0,\ldots,13$, satisfying
\begin{equation}
 \avg{F_A,F_B}=\eta_{AB},
 \qquad \eta_{AB}=\operatorname{diag}(-1,+1,\ldots,+1).
 \label{eq:profilebasis}
\end{equation}
\end{proposition}

\begin{proof}
Take a unit timelike constant profile for $F_0$.  For the positive directions
use normalized sine/cosine profiles in the three spacelike quaternion
basis directions, with distinct pairs (Fourier number, algebra direction).
Fourier orthogonality makes different choices orthogonal, while the Lorentz
slice identity fixes every norm to $+1$.  There are arbitrarily many such
positive profiles, so thirteen can be selected.
\end{proof}

Let $\iota:\mathbb R^{13,1}\to\mathscr V_\psi$ denote the isometric linear map
\[
 \iota(v):=v^A F_A,
 \qquad
 \avg{\iota(v),\iota(w)}=\eta_{AB}v^Aw^B.
\]

\section{Free-embedding lift}

\begin{definition}[Free isometric embedding]
A local isometric embedding $Y:U\to\mathbb E^{13,1}$ is free when the fourteen
ambient vectors
\begin{equation}
 \{Y_{,\mu},Y_{;\mu\nu}\},
 \qquad \mu=0,\ldots,3,\quad \mu\le\nu,
 \label{eq:freevectors}
\end{equation}
are linearly independent.  The ten $Y_{;\mu\nu}$ are normal to the four
$Y_{,\mu}$.
\end{definition}

\begin{theorem}[Profile lift of a free embedding]
Let $Y:U\to\mathbb E^{13,1}$ be a local free isometric embedding of a
Lorentzian metric $g$.  Define one biquaternionic profile field
\begin{equation}
 \boxed{\quad
 \Theta_Y(x,\psi):=\sqrt{\Nzero}\,Y^A(x)F_A(\psi).
 \quad}
 \label{eq:thetalift}
\end{equation}
Then:
\begin{enumerate}
 \item $g_{\mu\nu}[\Theta_Y]=g_{\mu\nu}$ exactly;
 \item the normal profile second derivatives are
 \begin{equation}
  B_{\mu\nu}
  :=P_\perp(\nabla_\mu E_\nu)
  =\sqrt{\Nzero}\,\iota(Y_{;\mu\nu});
  \label{eq:Btransfer}
 \end{equation}
 \item the closure map
 \begin{equation}
  \mathcal K(X)=X^{\mu\nu}B_{\mu\nu}
  \label{eq:Kmap}
 \end{equation}
 is injective.
\end{enumerate}
\end{theorem}

\begin{proof}
Differentiating Eq.~\eqref{eq:thetalift} gives
$E_\mu=\sqrt{\Nzero}\,\iota(Y_{,\mu})$.  The isometry of $\iota$ yields
\[
 \Nzero^{-1}\avg{E_\mu,E_\nu}
 =\eta_{AB}Y^A{}_{,\mu}Y^B{}_{,\nu}=g_{\mu\nu}.
\]
Because $\iota$ is constant and isometric, it commutes with spacetime
covariant differentiation and with orthogonal decomposition relative to the
tangent span.  Hence Eq.~\eqref{eq:Btransfer} holds.  Freeness makes the ten
normal vectors $Y_{;\mu\nu}$ linearly independent, and the injectivity of
$\iota$ preserves that independence.  Therefore $\mathcal K(X)=0$ implies
$X^{\mu\nu}=0$.
\end{proof}

\begin{theorem}[Local representability]
Every four-dimensional Lorentzian manifold of the regularity covered by the
local free-embedding theorem admits, near each point, a field of the form
Eq.~\eqref{eq:thetalift} that reproduces its metric and is fiber-free.
\end{theorem}

\begin{proof}
The local free-isometric-embedding theorem gives a free isometric embedding
$Y:U\to\mathbb E^{13,1}$.  Apply the profile-lift theorem.
\end{proof}

\begin{corollary}[Local Schwarzschild and general Einstein patches]
Every regular coordinate patch of Schwarzschild, Kerr, FLRW, or any other
smooth Lorentzian Einstein solution has a local fiber-free $\Theta_Y$
representer in the candidate profile metric.  This is an existence theorem;
it does not supply a preferred or global closed-form representer.
\end{corollary}

\section{Gauss identity and a direct pure-$\Theta$ action}

The target profile space has a fixed flat pairing.  Its Gauss equation gives
\begin{equation}
 R_{\mu\nu\rho\sigma}
 =\frac1{\Nzero}\left(
   \avg{B_{\mu\rho},B_{\nu\sigma}}
  -\avg{B_{\mu\sigma},B_{\nu\rho}}
 \right).
 \label{eq:Gauss}
\end{equation}
Consequently
\begin{equation}
 R[g[\Theta]]
 =\frac1{\Nzero}\left(
  \avg{B^\mu{}_{\mu},B^\nu{}_{\nu}}
  -\avg{B_{\mu\nu},B^{\mu\nu}}
 \right).
 \label{eq:scalarGauss}
\end{equation}

\begin{definition}[Pure-$\Theta$ Gauss action]
On the regular profile branch define
\begin{equation}
 \boxed{
 S_{\rm G}[\Theta]
 =\frac1{2\kappa}\int_U\sqrt{-g[\Theta]}\left[
 \frac1{\Nzero}\left(
  \avg{B^\mu{}_{\mu},B^\nu{}_{\nu}}
  -\avg{B_{\mu\nu},B^{\mu\nu}}
 \right)-2\Lambda\right]\dd^4x .}
 \label{eq:Gaussaction}
\end{equation}
The usual composite boundary term is understood, or variations are taken with
compact support.
\end{definition}

\begin{proposition}[No independent metric or connection]
On every regular profile configuration, Eq.~\eqref{eq:Gaussaction} equals the
Einstein--Hilbert--$\Lambda$ functional of the composite metric
$g[\Theta]$.  All objects in Eq.~\eqref{eq:Gaussaction} are built from one
$\Theta$, the fixed profile pairing, and spacetime derivatives.
\end{proposition}

\begin{theorem}[Exact local vacuum GR closure]
Let $\Theta$ be a regular fiber-free stationary point of
$S_{\rm G}[\Theta]$.  Then
\begin{equation}
 \boxed{G_{\mu\nu}+\Lambda g_{\mu\nu}=0.}
 \label{eq:Einstein}
\end{equation}
Conversely, every local free profile lift of an Einstein--$\Lambda$ metric is
a stationary point under compactly supported profile variations.
\end{theorem}

\begin{proof}
The metric variation is
\[
 \delta g_{\mu\nu}
 =\frac2{\Nzero}\avg{E_{(\mu},\partial_{\nu)}\delta\Theta}.
\]
Varying the composite Einstein--Hilbert functional, integrating by parts, and
using the contracted Bianchi identity gives the normal Euler--Lagrange
equation
\begin{equation}
 (G^{\mu\nu}+\Lambda g^{\mu\nu})B_{\mu\nu}=0.
 \label{eq:RTfree}
\end{equation}
Fiber-freeness makes the ten $B_{\mu\nu}$ independent, so
Eq.~\eqref{eq:RTfree} implies Eq.~\eqref{eq:Einstein}.  The converse follows
immediately because the Einstein residual vanishes before contraction with
$B_{\mu\nu}$.
\end{proof}

\section{What is closed and what remains}

\begin{center}
\begin{tabularx}{\textwidth}{>{\raggedright\arraybackslash}p{0.28\textwidth}
                              >{\raggedright\arraybackslash}p{0.20\textwidth}
                              >{\raggedright\arraybackslash}X}
\hline
Item & Status & Scope \\
\hline
Local metric representability & Closed (profile metric) & Follows from local free embedding in $\mathbb E^{13,1}$ and the profile lift. \\
Fiber-free non-vacuity on GR geometries & Closed locally & Every locally free embedded GR patch lifts to one $\Theta$. \\
Vacuum Einstein--$\Lambda$ equation & Closed (free branch) & Direct variation of the pure-$\Theta$ Gauss action. \\
Independent tetrad/connection fields & None introduced & $g$, $B$, and the Levi--Civita objects are composites of $\Theta$. \\
Canonical metric choice & Axiom decision & Current AXIOM C explicitly forbids compact-$\psi$ averaging. \\
Origin of $\kappa$ and $\Lambda$ & Open & The Gauss action fixes the form, not their values from the older master equation. \\
Dynamic profile selection/stability & Open & Must show the UBT dynamics selects and preserves the free, $\psi$-stable branch. \\
Holomorphic/Jacobi restriction & Open & The lift is smooth periodic; any stronger analytic restriction must be checked separately. \\
Global representability & Open & The theorem is local; topology and horizons require patching/global analysis. \\
Matter and gauge separation & Open & Vacuum needs no separate $\mathcal F_\Theta$ term; unified matter requires a projection/separation theorem. \\
\hline
\end{tabularx}
\end{center}

\section{Decision}

The candidate profile completion is not merely a right inverse for a prescribed
tetrad.  Once the profile metric and the Gauss action are adopted, the field
$\Theta$ is the embedding variable, its stationary equation is
Eq.~\eqref{eq:RTfree}, and freeness converts that equation to the complete ten
Einstein equations.  The construction therefore supplies an exact local
vacuum GR sector from one $\Theta$.

The remaining obstacle is conceptual and canonical rather than a missing
rank calculation: decide whether the classical metric is a pointwise tetrad
readout or the translation-invariant readout of the full already-existing
complex-time profile.  The former has the proved torsion-free concurrent-vector
no-go; the latter has the exact local closure proved above.

\begin{thebibliography}{9}
\bibitem{Bustamante2005}
M. D. Bustamante, P. Debbasch, and M. E. Brachet,
``Classical Gravitation as Free Membrane Dynamics,''
\emph{Class. Quantum Grav.} 23 (2006), arXiv:gr-qc/0509090.

\bibitem{Willison2013}
S. Willison,
``A Re-examination of the Isometric Embedding Approach to General Relativity,''
arXiv:1311.6203.
\end{thebibliography}

\end{document}
